\chapter{化学基础知识}

\section{气体}

% 气体基本特性：
% \begin{enumerate}
%     \item 自动扩散性
%     \item 可压缩性
% \end{enumerate}

\subsection{气体的状态方程}
\subsubsection{理想气体的状态方程}
假设分子微粒之间的作用力、分子本身的体积忽略不计的气体为\textbf{理想气体}。

（\textbf{低压高温}状态）理想气体状态方程 $$\boxed{pV=nRT}$$

\subsubsection{实际气体的状态方程}
Van der Waals 气体状态方程 $$\boxed{\left(p+a\left({\dfrac{n}{V}}\right)^2\right)(V-bn)=nRT}$$其中 $p_{\text{内}}=\left(\dfrac{n_{\text{外}}}{V}\right)\left(\dfrac{n_{\text{内}}}{V}\right)=\left(\dfrac{n}{V}\right)^2$，
$a$ 是与分子间引力有关的常数，$b$ 是与分子本身体积有关的常数，统称为 van der Waals 常数。通常，$a$ 和 $b$ 数值越大，表明实际气体距理想气体愈远，从而产生的偏差就愈大。

\subsection{混合气体的分压定律}
Dalton 分压定律 $$\boxed{p_{\text{总}}=\sum_i p_i}$$即$$\boxed{p_i=p_{\text{总}}x_i}$$

Amagat 分体积定律 $$\boxed{V_{\text{总}}=\sum_i V_i}$$

\subsection{气体扩散定律}
$$\boxed{v \propto \sqrt{\dfrac{1}{\rho}}}$$
$$\boxed{v \propto \sqrt{\dfrac{1}{M}}}$$

\subsection{气体分子的速率分布和能量分布}
\subsubsection{气体分子的速率分布}
1859 年，英国物理学家 Maxwell 用概率论的方法推导出了计算气体分子运动速率分布的公式，讨论了分子运动速率的分布。后来，奥地利人 L. Boltzmann 用统计力学的方法也得到了相同的结果，从而进一步验证了 Maxwell 的速率分布公式。

速率分布函数满足归一化条件：$$\int_{0}^{\infty} f(v) {\rm d}v = 1$$

\begin{itemize}
    \item 最概然速率：$v_p = 1.41\sqrt{\dfrac{RT}{M}}$
    \item 算术平均速率：$\bar{v} = 1.60\sqrt{\dfrac{RT}{M}}$
    \item 方均根速率：$\sqrt{\overline{v^2}} = 1.73\sqrt{\dfrac{RT}{M}}$
\end{itemize}

比例关系：$$\sqrt{\overline{v^2}}:\bar{v}:v_p = 1.73:1.60:1.41$$

\subsubsection{气体分子的能量分布}
在无机化学中，常用下面的近似公式来计算和讨论能量的分布。
\[
f_{E_0}=\dfrac{N_i}{N}=\exp{\left(\dfrac{-E_0}{RT}\right)}
\]

\section{液体和溶液}

\subsection{相与态}

相：系统中物理性质、化学性质完全相同部分，相间有界面

态：物质的聚集状态

\subsection{溶液浓度的表示方法}

\textbf{溶液}是指一种或一种以上的物质以分子或离子的形式溶解在另一种物质中形成的均一、稳定的混合物。其中物质的量多且能溶解其他物质的组分称为\textbf{溶剂}，含量少的组分称为\textbf{溶质}。

对于溶剂 $\mathrm{A}$ 和溶质 $\mathrm{B}$，有
\begin{enumerate}
    \item 质量摩尔浓度 $\boxed{b({\rm B})=\dfrac{n({\rm B})}{m({\rm A})}}$
    \item 物质的量浓度 $\boxed{c({\rm B})=\dfrac{n({\rm B})}{V}}$
    \item 质量分数 $\boxed{w({\rm B})=\dfrac{m({\rm B})}{m}}$
    \item 摩尔分数 $\boxed{x({\rm B})=\dfrac{n({\rm B})}{n}}$
\end{enumerate}
\[
x \approx \dfrac{1000}{18}b~\mathrm{mol^{-1}}
\]

\subsection{溶液的饱和蒸气压}
\subsubsection{纯溶剂的饱和蒸气压}
当凝聚速率和蒸发速率相等时，饱和蒸气所具有的压强叫做该温度下液体的\textbf{饱和蒸气压}，简称\textbf{蒸气压}，用 $p^*$ 表示。对同一液体来说，其饱和蒸气压随温度的升高而增大。

\subsubsection{溶液的饱和蒸气压下降}
溶液的饱和蒸气压 $p$ 小于纯溶剂的蒸气压 $p^*$。溶液的浓度越大，溶质的粒子数目越多，溶液的饱和蒸气压比纯溶剂的饱和蒸气压下降得就越多。

Raoult 定律 $\boxed{p=p^*x({\rm A})}$

\[
\begin{aligned}
\Delta p &= p^*-p \\
&=p^*[1-x({\rm A})] \\
&=p^*x({\rm B})
\end{aligned}
\]

稀溶液的饱和蒸气压下降值与稀溶液的质量摩尔浓度成正比。

\subsection{稀溶液的依数性}
稀溶液的饱和蒸气压下降是稀溶液的一种依数性质。

\subsubsection{溶液的沸点升高}
液体汽化的方式有两种，即\textbf{蒸发}和\textbf{沸腾}。当液体的饱和蒸气压与外界大气压强相等时，液体的汽化将在其表面和内部同时发生，称为\textbf{沸腾}，这时的温度称为该液体的\textbf{沸点}。而在低于此温度下的汽化，仅在液体表面上进行，称为\textbf{蒸发}。

难挥发性非电解质稀溶液沸点升高值 $\Delta T_{\mathrm{b}}=T_{\mathrm{b}}-T_{\mathrm{b}}^*$，与其饱和蒸气压下降的数值成正比，即 $\Delta T_{\mathrm{b}} \propto \Delta p$，记 $k_{\mathrm{b}}$ 为稀溶液\textbf{沸点升高常数}，得\textbf{沸点升高公式} $$\boxed{\Delta T_{\mathrm{b}}=k_{\mathrm{b}}b}$$

\subsubsection{溶液的凝固点降低}
液体凝固成固体（严格说是晶体）时的温度称为该液体的\textbf{凝固点}。由于溶液的蒸气压下降导致溶液的凝固点低于纯溶剂的凝固点。

难挥发性非电解质稀溶液凝固点降低值 $\Delta T_{\mathrm{f}}=T_{\mathrm{f}}^*-T_{\mathrm{f}}$，与其饱和蒸气压下降的数值成正比，记 $k_{\mathrm{f}}$ 为\textbf{凝固点降低常数}，得\textbf{凝固点降低公式} $$\boxed{\Delta T_{\mathrm{f}}=k_{\mathrm{f}}b}$$

注：$k_{\mathrm{b}}$、$k_{\mathrm{f}}$ 仅与溶剂本性有关，与溶质种类无关；$\Delta T_{\mathrm{b}}$、$\Delta T_{\mathrm{f}}$ 与溶质粒子总数有关，可用于测定难挥发非电解质分子量。

\subsubsection{渗透压}
溶剂分子透过半透膜进入溶液的现象，称为\textbf{渗透现象}。渗透平衡时两侧液面高度差造成的静压称为溶液的\textbf{渗透压}。

Van't Hoff 指出 $$\boxed{\varPi V = nRT}$$ 即 $$\boxed{\varPi = c_{\rm B}RT}$$

渗透压对浓度敏感，低浓度即可产生显著渗透压，如 $c_{\rm B}=0.00100\ \mathrm{mol·L^{-1}}$ 时，$\varPi=2478\ \mathrm{Pa}$。

\section{固体和晶体}
\subsection{晶体和非晶体}
\subsubsection{非晶体的特性}
原子排列紊乱分布，无明显固定熔点，受热时逐渐软化；物理性质表现为各向同性。

\subsubsection{晶体的特性}
原子呈规则周期性排列，有固定熔点；物理性质（如光学、力学性质）具有各向异性。

\subsection{对称性}
\subsubsection{旋转和旋转轴}
旋转是晶体常见对称操作，绕旋转轴旋转一定角度后晶体复原。宏观对称旋转轴仅存在1、2、3、4、6次轴，无5次或更高次轴。

\subsubsection{反映和镜面}
镜面（对称面，记为$m$）是将晶体分为互为镜像的两部分的平面，通过镜面的反映操作后晶体复原。

\subsubsection{反演和对称中心}
对称中心（记为$i$）是晶体内部一点，通过该点的反演操作（任一点沿连线过中心到另一侧等距处）后晶体复原。

\subsection{晶体和点阵}
晶体的规则结构可抽象为\textbf{点阵}，点阵由一系列等同位置（阵点）组成，阵点代表晶体中原子、分子等的等同排列位置。

\subsection{晶系和点阵型式}
\subsubsection{七个晶系}
按晶格参数（晶胞边长 $a$、$b$、$c$ 和夹角 $\alpha$、$\beta$、$\gamma$）分为七类：

\begin{table}[h!]
  \centering
  \caption*{晶系与晶格参数关系}
  \begin{tabular}{lc}
    \toprule
    \textbf{晶系} & \textbf{晶格参数关系} \\
    \midrule
    三斜 & $a \neq b \neq c$，$\alpha \neq \beta \neq \gamma \neq \ang{90}$ \\
    单斜 & $a \neq b \neq c$，$\alpha = \gamma = \ang{90} \neq \beta$ \\
    正交 & $a \neq b \neq c$，$\alpha = \beta = \gamma = \ang{90}$ \\
    四方 & $a = b \neq c$，$\alpha = \beta = \gamma = \ang{90}$ \\
    立方 & $a = b = c$，$\alpha = \beta = \gamma = \ang{90}$ \\
    三方 & $a = b = c$，$\alpha = \beta = \gamma \neq \ang{90}$ \\
    六方 & $a = b \neq c$，$\alpha = \beta = \ang{90}$，$\gamma = \ang{120}$ \\
    \bottomrule
  \end{tabular}
\end{table}

\subsubsection{十四种空间点阵型式}
七个晶系对应14种空间点阵（Bravais点阵）：
\begin{itemize}
    \item 三斜：
        \begin{itemize}
            \item 简单三斜
        \end{itemize}
    \item 单斜：
        \begin{itemize}
            \item 简单单斜
            \item 底心单斜
        \end{itemize}
    \item 正交：
        \begin{itemize}
            \item 简单正交
            \item 体心正交
            \item 底心正交
            \item 面心正交
        \end{itemize}
    \item 四方：
        \begin{itemize}
            \item 简单四方
            \item 体心四方
        \end{itemize}
    \item 立方：
        \begin{itemize}
            \item 简单立方
            \item 体心立方
            \item 面心立方
        \end{itemize}
    \item 三方：
        \begin{itemize}
            \item 简单三方
        \end{itemize}
    \item 六方：
        \begin{itemize}
            \item 简单六方
        \end{itemize}
\end{itemize}

\subsection{晶胞}
\textbf{晶胞}是点阵的基本重复单元，能反映点阵的全部特征。

晶胞参数：即晶格参数 $a$、$b$、$c$（边长）和 $\alpha$、$\beta$、$\gamma$（夹角）。

晶胞类型：
\begin{description}
    \item[素晶胞] 只含1个阵点的晶胞；
    \item[复晶胞] 含2个及以上阵点的晶胞（如体心、面心晶胞）。
\end{description}

选取原则：反映点阵对称性、体积最小、棱边夹角接近 \ang{90}。